\section{Conclusion} This study explored the promise of artificial
intelligence/machine learning techniques to identify optimal designs for a
single channel of a \gls{MSR}. We first created a reactor database to train and
test \gls{ML} methods and searched for the best performing estimator among those
\gls{ML} methods by evaluating their scores, cross-validation plots, and
confusion matrices. After deciding on the estimator, we used the genetic
algorithm technique (\gls{NSGAIII}) coupled to the selected estimator (\gls{ET})
to seek for the optimum single channel designs. We found a set of optimum
solutions for each fuel-salt pair considering several design constraints on the
performance metrics (i.e., fuel type, salt type, $^{235}$U fraction in U, U
fraction in heavy metal, channel-to-channel pitch length, moderator-to-salt
(volume) ratio, and average fuel-salt temperature). We accurately classified and
predicted design performance metrics. Major findings throughout this study and
general remarks on the results follow:

\paragraph{ As a result of the assessment of the scores of the estimators,
	\gls{RF}, \gls{ET}, \gls{KN}, \gls{MLP}, and \gls{SV} were chosen as the best
	candidate estimators for the design optimization phase} Estimators of ensemble,
support vector, \gls{NN}, and decision tree methods performed well, while the
linear methods (L, R, SGD etc.) performed poorly as compared to the non-linear
methods (\gls{DT}, GB, and \gls{RF} etc.). \paragraph{We evaluated a total
	optimum sample size of 300,000 for sufficiently accurate results} The learning
curves of the selected estimators suggested an optimum sample size in the
dataset of at least 12,000 for each fuel-salt pair. \paragraph{The \gls{ET}
	estimator was selected as the best estimator} The cross-validation plots and
confusion matrices of the candidate estimators shed light on the best estimator
for the design optimization phase by pinpointing the individual performance
metric. All the performance metrics were predicted with less than 5\% error and
classified with less than 1\% error except for the feedback label which has an
error of 10\%.  We also understood that the \gls{MLP} handles performance
metrics as good as \gls{ET}, but with considerably higher computational
run-time. \paragraph{We estimated all the performance metrics of the optimum
	designs within a prediction error of at most 5\% in comparison to their actual
	values} A comprehensive overview of the common properties of the optimal designs
implied a particular enrichment level and U content associated with fuel type
but not salt type, $T_{ave}$ higher than 1100 K, and a moderator-to-salt ratio
of 0.274. The relation of the channel-to-channel pitch length, $p$, with either
fuel type or salt type was not as clear as the other feature variables.
\paragraph{All the optimum designs were accurately labeled as supercritical
	reactor, burner type, fast spectrum, and negative flag in feedback coefficient}
As indicated by confusion matrices, only a few, most of which belong to the
feedback coefficient, were mislabeled.

\paragraph{Based on the constraints in this study, the optimal fuel-salt option
	was a combination of U-Pu fuel and NaCl salt} It achieved the longest graphite
lifetime (i.e., the lowest $\phi_{fast}$), the highest $\varsigma$, and a
sufficient negative ${\alpha_{total}}$.

In addition to the above numerical findings, the most important success of this
study is the integration of the \gls{ML}-enabled technique into the reactor core
design process. Another important success achieved by this proposed method is
the examination of the optimal designs in a very short time (within minutes),
regardless of the value and number of performance metrics or feature parameters,
once a reactor database is generated. Without using any reactor physics code,
this method makes available evaluating the performance metrics of any design
with very high accuracy by using a training dataset specific to the feature
parameters of interest to be loaded from a reactor database generated initially.
Due to the very fast computing speed and high accuracy in estimating performance
metrics, we recommend using this method for neutronic calculations (also for
thermal-hydraulics) in nuclear engineering along with nuclear computer codes
such as deterministic and Monte Carlo that require a significant user expertise.

In addition to these achievements, we encountered various difficulties that have
an direct and indirect impact on the outcomes of this study throughout the work.
First, we understood that the number of performance metrics, sample size, and
assigning different estimators (composite estimators) to an individual
performance metric significantly affect the errors of performance metrics. For
instance, the larger the sample size in the training dataset, the higher the fit
scores, thus better fitting results. Second, there were several known drawbacks
to getting better estimators' scores: one was the hyper-parameter dependency
influencing the scores of the estimators and the other one was the dependency on
the quality of training datasets containing out-of-interest values (outliers),
such as very low $k_{\infty}$ and very high $\varsigma$. Although the effects of
hyper-parameters are limited, outliers have a profound effect on the fit scores,
therefore on the predicted performance metrics and optimal designs.

We also figured out throughout this study that the primary difficulty with
regard to the applicability of \gls{ML} methods for reactor design is to
generate a reactor database which consumes the entire computer's resources,
requires a powerful supercomputer, and takes a long time. In fact, this is not a
critical issue for the applicability of the proposed method since the database
generation is quite easy process in comparison to other tasks and it is enough
to do it once at the very beginning of the study. Third, it is likely to find
different solutions when the population size and number of generations are
further increased, or a different optimization method, such as brute force which
scans the entire grid space for a global optimum solution, is used. But, we
anticipate that all the solutions to be obtained using different parameters
would be close to the global optima presented in this work as we already
verified our solutions with different optimizers and their varying
hyper-parameters. The only concern with the optimization process is the defined
constraints on performance metrics as an optimizer directs offspring according
to these constraints.

In our follow-up studies, we will take the following steps to do away with some
of the assumptions and simplifications we have made in this research, to
complement the shortcomings in the study and to further improve the existing
method. Initially, we have a plan to use the suggested method for a full-core
(or small scale or assembly) \gls{MSR} design in the second phase of the ongoing
study. Different from the single channel design, designing a full reactor core
will bring several additional variables (e.g., number of fuel channels, reactor
diameter) to the existing feature parameters and a few additional performance
metrics (e.g., power peaking factor) to the existing performance metrics. The
rest is expected to follow the method suggested in this study. Concurrently, we
will replace the grid-based search technique used in generating the reactor
database by a more sophisticated sampling strategy (e.g., adaptive, stratified
and hybrid), which would help reduce the number of datasets, by focusing only on
a particular domain. This is because, despite being the simplest and capable of
evaluating all the possible scenarios among the different parameters of feature
space, the technique is computationally expensive and consumes computing
resources by a considerable amount. On the other hand, in order to simplify the
optimization problem and make a significant gain in the simulation run-time, we
had deliberately disregarded the thermal-hydraulics effects such as delayed
precursor movement, temperature distribution along the channel, and reactivity
feedback on the neutronic calculations. Related to this, we will integrate the
Moltres MOOSE application that solves the thermal-hydraulics and neutronics
equations for liquid-fueled \gls{MSR}s with the Plankton python tool having been
developed in this work and increase the accuracy of neutronic calculations.

To sum up, this study has explored an approach to designing a nuclear reactor
core from scratch by incorporating the \gls{AI}/\gls{ML} techniques into single
channel analysis. The main focus of this study was to compare the prediction
accuracy of different \gls{ML} methods with each other and to decide the best
estimator suitable for single channel analysis. The research will make a
significant contribution to the further advancement of optimization studies by
eliminating the deficiencies of existing methods in reaching ideal designs such
as the tournament approach of the optimizers, convergence criteria, strict
constraints on performance metrics and concerns about computer resource
utilization. The results of this study could benefit researchers who are
interested in a multi-purpose reactor design which provides more flexibility for
in-core fuel management, improved neutron economy for fuel utilization, more
safety margin against core melt-down, prolonged reactor life-time, and less
spent nuclear fuel.
